\documentclass[10pt,a4paper]{article}

\usepackage{cmap}
\usepackage[left=1.4cm,top=0.7cm,right=1.4cm,bottom=0.45cm]{geometry}
\usepackage[utf8]{inputenc}
\usepackage[T1]{fontenc}
\usepackage[spanish]{babel}
\usepackage{hyperref}
\usepackage{enumitem}
\usepackage{titlesec}

\setlength{\parindent}{0pt}
\setlength{\parskip}{0pt}
\linespread{0.90}
\hypersetup{ 
    colorlinks=true,
    urlcolor=blue,
    linkcolor=blue
}

% ---- Formato de secciones ----
\titleformat{\section}{\large\bfseries}{}{0em}{}[\vspace{0.5pt}\hrule]
\titlespacing{\section}{0pt}{4pt}{2pt}

% ---- Comando para encabezados de experiencia/proyectos (texto plano, sin tablas, ATS-friendly) ----
\newcommand{\entryheader}[2]{\textbf{#1}\\ \textit{#2}\\[1pt]}

\begin{document}
\pagestyle{empty}

% ================= HEADER =================
\begin{center}
    {\LARGE \textbf{Giuliano Pascarelli}}\\
    \vspace{2pt}
    {\normalsize Desarrollador Full Stack}\\
    \vspace{1.5pt}
    \small
    San Luis, Argentina \textbullet\ giulianopascarelli@gmail.com \textbullet\ +54 2664 857382\\
    \href{https://lelox028.github.io/}{Portfolio} \textbullet\
    \href{https://www.linkedin.com/in/giuliano-pascarelli-097100190/}{LinkedIn} \textbullet\
    \href{https://github.com/lelox028}{GitHub}
\end{center}

\vspace{2pt}
\hrule

% ================= PERFIL =================
\section*{Perfil Profesional}
Desarrollador Full Stack con más de 3 años de experiencia construyendo aplicaciones web, con especial fortaleza en Frontend (React, Next.js) y bases sólidas en Backend (Java, Node.js, APIs REST). Experiencia liderando equipos pequeños de desarrollo, capacitando a más de 20 desarrolladores junior y dictando cursos de programación. Cómodo trabajando con bases de datos SQL, metodologías ágiles (Scrum/Kanban) y contribuyendo a proyectos de código abierto orientados a la comunidad.

% ================= EXPERIENCIA =================
\section*{Experiencia}

\entryheader{Desarrollador Backend / Full Stack --- SimpleDevsSL (Iniciativa Open Source)}{Diciembre 2025 -- Actualidad}
Iniciativa colaborativa y sin fines de lucro que desarrolla soluciones de software open source para problemáticas reales de la comunidad de San Luis, Argentina.
\begin{itemize}[noitemsep, topsep=1pt, partopsep=0pt, leftmargin=16pt]
    \item Colaboración con un equipo reducido de desarrolladores en el diseño y desarrollo de proyectos de software libre orientados a resolver necesidades reales de la comunidad.
    \item Contribución principal en desarrollo backend, arquitectura de software y planificación técnica, participando también en frontend, infraestructura y despliegue según necesidad del proyecto.
    \item Desarrollo colaborativo de Guía Puntana, plataforma web para conectar talento y servicios locales.
\end{itemize}

\vspace{2pt}
\entryheader{Desarrollador Full Stack / Team Leader --- Aptugo, Warwick, Reino Unido (remoto desde Argentina)}{Febrero 2021 -- Abril 2024}
Me incorporé a Aptugo con conocimientos iniciales de programación y, a través de un proceso de aprendizaje autodidacta y trabajo constante, crecí hasta liderar un equipo de desarrollo como Team Leader, combinando roles de desarrollo, liderazgo técnico y capacitación en un entorno startup de ritmo acelerado.
\begin{itemize}[noitemsep, topsep=1pt, partopsep=0pt, leftmargin=16pt]
    \item Desarrollo de aplicaciones web full stack con el stack MERN (MongoDB, Express, React, Node.js) e interfaces responsive con React, SCSS y Material UI.
    \item Diseño e integración de APIs REST y mantenimiento de bases de datos NoSQL (MongoDB).
    \item Coordinación técnica de un equipo de hasta 3 desarrolladores como Team Leader, y capacitación de más de 20 desarrolladores junior con niveles de experiencia previa muy variados.
    \item Dictado de cursos de programación a clientes externos sobre la herramienta low-code de la empresa, enseñando fundamentos de programación tradicional y su equivalencia con conceptos low-code.
    \item Trabajo bajo metodologías ágiles (Scrum y Kanban).
\end{itemize}

\vspace{2pt}
\entryheader{Profesor Particular de Programación Orientada a Objetos (POO) --- Independiente}{Freelance}
Clases particulares de POO a estudiantes y autodidactas, aplicando la experiencia adquirida como tutor y capacitador en Aptugo.

% ================= PROYECTOS =================
\section*{Proyectos}

\entryheader{Gym Tracker --- React, Vite, Hono, Cloudflare Workers}{\href{https://gym-tracker-frontend.giulianopascarelli.workers.dev/}{gym-tracker-frontend.giulianopascarelli.workers.dev}}
Aplicación web mobile-first para gestionar ejercicios y entrenamientos de gimnasio. Frontend en React y Vite, API en Hono con base de datos SQL (D1), con despliegue completo en Cloudflare Workers. Proyecto personal.

\vspace{2pt}
\entryheader{Study Plan Manager --- React, Node.js, Express}{\href{https://github.com/lelox028/study-plan-manager-frontend-react}{Repo Frontend} \textbullet\ \href{https://github.com/lelox028/study-plan-manager-api}{Repo Backend}}
Aplicación full stack personal para ayudar a estudiantes a organizar su plan de estudios y cronograma académico. Proyecto individual, actualmente en un 80\% de avance.

\vspace{2pt}
\entryheader{Juliette's Kitchen Timer --- React, Vite}{\href{https://julietteskitchentimer.vercel.app/}{julietteskitchentimer.vercel.app}}
PWA liviana basada en la técnica Pomodoro, con una estética ilustrada propia y una mascota animada que acompaña las sesiones de trabajo. Proyecto personal.

% ================= EDUCACIÓN =================
\section*{Educación}
\textbf{Licenciatura en Informática} (en curso) --- Universidad Siglo 21, Córdoba, Argentina\\
\textit{2024 -- Actualidad}
\begin{itemize}
    \item Título intermedio de \textbf{Analista en Sistemas} obtenido\\
    \textit{Julio 2026}
\end{itemize}

% \vspace{2pt}
% \textbf{Título Secundario} --- Colegio Don Bosco, San Luis, Argentina\\
% \textit{2012 -- 2018}

% ================= PERFIL TÉCNICO =================
\section*{Perfil Técnico}

\textbf{Lenguajes:} Java, JavaScript, TypeScript, C\\
\textbf{Frontend:} React, Next.js, HTML, CSS, SCSS, Material UI\\
\textbf{Backend:} Node.js, Express, Hono, Java (Spring Boot), APIs REST\\
\textbf{Bases de Datos:} SQL (MySQL), Supabase, MongoDB (nivel básico)\\
\textbf{Herramientas y metodologías:} Git, Docker (entorno personal), Cloudflare Workers, Maven, UNIX/Linux, shell scripting, Scrum, Kanban

\vspace{1.5pt}
\textbf{Habilidades blandas:} Liderazgo técnico y coordinación de equipos, mentoría y capacitación de perfiles junior, comunicación efectiva y trabajo en equipo, capacidad de aprendizaje autodidacta y resolución de problemas.

\vspace{1.5pt}
\textbf{Idiomas:} Español (Nativo), Inglés (Avanzado -- lectura, escritura y conversación)

\vspace{1.5pt}
\textbf{Intereses:} Proyectos de código abierto, aprendizaje continuo de nuevas tecnologías y personalización de sistemas Linux.

\end{document}
